%============================================================
% CHƯƠNG 9: ĐẠO ĐỐI NHÂN XỬ THẾ
%============================================================
\chapter{Đạo Đối Nhân Xử Thế (The Way of Dealing with People)}

\begin{center}
  \textit{\color{truyenblue}``You can easily judge the character of a man by how he treats those who can do nothing for him.''}\\
  \textit{(Bạn có thể dễ dàng đánh giá nhân cách của một người qua cách họ đối xử với những ai không thể làm gì cho họ.)}\\[4pt]
  \small\color{truyengold}--- Malcolm Forbes ---
\end{center}
\ngancach

\section{Người Bồi Bàn}
\begin{truyen}{Người Bồi Bàn}{Nhân Cách}
\chuhoa{T}{ }rước khi quyết định kết hôn với một người, người cha chỉ yêu cầu quan sát anh ta một lần ăn tối.\\
\textit{(Before agreeing to his daughter's marriage, the father only asked to observe the suitor once during dinner.)}

Trong bữa ăn đó, ông không nhìn con gái mà nhìn cách chàng trai đối xử với nhân viên bồi bàn.\\
\textit{(Throughout the meal, he watched not his daughter but how the young man treated the waitstaff.)}

Chàng trai courtesy, nói cảm ơn, nhẹ nhàng khi có nhầm lẫn.\\
\textit{(The young man was courteous, said thank you, and was gracious when there were mistakes.)}

``\textbf{Cách anh ta đối xử với người phục vụ mới là cách anh ta sẽ đối xử với con gái tôi khi không còn cần gây ấn tượng nữa},'' ông giải thích.\\
\textit{(``\textbf{How he treats the server is how he will treat my daughter when he no longer needs to impress anyone},'' he explained.)}
\end{truyen}
\begin{baihoc}
  Nhân cách thực sự của con người bộc lộ rõ nhất trong cách họ đối xử với những người không có quyền lực hay ảnh hưởng với họ. Đó là thước đo đáng tin cậy nhất.
\end{baihoc}

\section{Lời Khen Chân Thật}
\begin{truyen}{Lời Khen Chân Thật}{Chân Thành}
\chuhoa{H}{ }ai đồng nghiệp cùng nộp đề xuất cho dự án. Sếp chọn đề xuất của người kia.\\
\textit{(Two colleagues both submitted proposals for a project. The boss chose the other person's.)}

Người không được chọn tìm đến người được chọn và nói thật lòng: ``Đề xuất của cậu tốt hơn của tôi. Tôi đặc biệt thích cách cậu xử lý phần ngân sách.''\\
\textit{(The one not chosen sought out the one who was and said sincerely: ``Your proposal was better than mine. I especially liked how you handled the budget section.'')}

Người kia ngạc nhiên: ``Thật không?''\\
\textit{(The other was surprised: ``Really?'')}

``Tôi nói thật. Tôi thất vọng vì không được chọn, nhưng thất vọng không có nghĩa là tôi phải giả vờ rằng cậu không xứng đáng.''\\
\textit{(``I am being honest. I am disappointed for not being chosen, but disappointment does not mean I should pretend you did not deserve it.'')}

Từ đó hai người trở thành \textbf{đồng minh suốt sự nghiệp}.\\
\textit{(From then on the two became \textbf{allies throughout their careers}.)}
\end{truyen}
\begin{baihoc}
  Lời khen chân thành trước người hơn mình đòi hỏi dũng cảm và phẩm cách. Sự lớn lượng trong thất bại là một trong những biểu hiện đẹp nhất của nhân cách.
\end{baihoc}

\section{Tin Nhắn Sai Địa Chỉ}
\begin{truyen}{Tin Nhắn Sai Địa Chỉ}{Tử Tế Với Người Lạ}
\chuhoa{C}{ }ô nhận tin nhắn từ số lạ: ``Anh ơi, em đang khóc một mình không biết phải làm gì.''\\
\textit{(She received a text from an unknown number: ``Brother, I am crying alone and do not know what to do.'')}

Rõ ràng là nhắn nhầm số.\\
\textit{(It was clearly sent to the wrong number.)}

Cô có thể bỏ qua. Nhưng cô nhắn lại: ``Chắc bạn nhắn nhầm số rồi. Nhưng tôi ở đây nếu bạn muốn kể.''\\
\textit{(She could have ignored it. But she replied: ``You probably sent this to the wrong number. But I am here if you want to talk.'')}

Họ nhắn qua lại một tiếng.\\
\textit{(They texted back and forth for an hour.)}

Người kia cuối cùng nhắn: ``Cảm ơn. Nhắn nhầm mà gặp đúng người cần gặp.''\\
\textit{(The other person finally wrote: ``Thank you. Wrong number but the right person.'')}
\end{truyen}
\begin{baihoc}
  Tử tế với người lạ không cần lý do hay mối liên hệ. Đôi khi một người xa lạ ở đúng nơi đúng lúc có thể cứu một ngày, thậm chí cứu một cuộc đời.
\end{baihoc}

\section{Người Hỏi Đường}
\begin{truyen}{Người Hỏi Đường}{Kiên Nhẫn}
\chuhoa{M}{ }ột khách du lịch người nước ngoài hỏi đường theo tiếng Việt khó hiểu.\\
\textit{(A foreign tourist asked for directions in broken Vietnamese.)}

Người đầu tiên được hỏi lắc đầu đi thẳng. Người thứ hai cũng vậy.\\
\textit{(The first person shook their head and walked on. So did the second.)}

Người thứ ba dừng lại và dù không nói được tiếng Anh, \textbf{mở điện thoại Google Maps ra chỉ từng bước}.\\
\textit{(The third person stopped and despite not speaking English, \textbf{opened Google Maps and pointed out each step}.)}

Họ đứng ở góc đường mười lăm phút, trao đổi bằng bản đồ và cử chỉ.\\
\textit{(They stood at the corner for fifteen minutes, communicating through maps and gestures.)}

Khách nước ngoài cúi đầu cảm ơn nhiều lần. Người chỉ đường gật đầu mỉm cười, \textbf{cảm thấy vui hơn cả ngày hôm đó}.\\
\textit{(The tourist bowed in thanks many times. The helper nodded and smiled, \textbf{feeling happier than at any other moment that day}.)}
\end{truyen}
\begin{baihoc}
  Sự kiên nhẫn và thiện chí với người lạ không chỉ giúp người được nhận mà còn mang lại niềm vui thực sự cho người cho đi. Lòng tốt tự nó là bản thưởng.
\end{baihoc}

\section{Lời Xin Lỗi Công Khai}
\begin{truyen}{Lời Xin Lỗi Công Khai}{Dũng Cảm}
\chuhoa{T}{ }rong buổi họp toàn công ty, người quản lý đứng dậy xin lỗi nhân viên vì đã đánh giá sai và công khai khiển trách người đó tuần trước.\\
\textit{(During the company-wide meeting, the manager stood up to apologise to an employee for having misjudged and publicly reprimanded that person the previous week.)}

Cả phòng im lặng bất ngờ.\\
\textit{(The whole room fell into surprised silence.)}

``Tôi sai. Tôi xin lỗi trước mọi người vì tôi đã khiến anh ấy xấu hổ trước mọi người.''\\
\textit{(``I was wrong. I apologise before everyone because I shamed him before everyone.'')}

Không ai ngờ một sếp cấp cao lại làm vậy. Nhưng từ hôm đó, \textbf{nhân viên tin tưởng vào công ty hơn bất kỳ package phúc lợi nào}.\\
\textit{(No one expected a senior manager to do that. But from that day, \textbf{employees trusted the company more than any benefit package could buy}.)}
\end{truyen}
\begin{baihoc}
  Xin lỗi công khai khi phạm lỗi công khai đòi hỏi dũng cảm lớn hơn bất kỳ quyết định kinh doanh nào. Người dám thừa nhận sai lầm là người thực sự mạnh mẽ.
\end{baihoc}

\section{Người Bán Hàng Rong}
\begin{truyen}{Người Bán Hàng Rong}{Tôn Trọng}
\chuhoa{M}{ }ột cô gái trẻ dừng xe mua bánh của bà cụ bán rong.\\
\textit{(A young woman stopped her car to buy pastries from an elderly street vendor.)}

Cô không vội vã lấy hàng rồi đi mà \textbf{dừng lại nói chuyện vài phút hỏi thăm bà}.\\
\textit{(She did not rush to grab the goods and leave but \textbf{stopped to chat for a few minutes and ask after the grandmother}.)}

Khi cô đi rồi, bà cụ về kể với con gái: ``Hôm nay có cô khách lịch sự quá, hỏi thăm bà như hỏi thăm người thân vậy.''\\
\textit{(After the young woman left, the grandmother told her daughter: ``Today there was such a polite customer, she asked after me as if I were family.'')}

``Mà cô ấy có mua nhiều không hả mẹ?'' – ``Chỉ mua một cái bánh thôi. Nhưng bà vui cả buổi.''\\
\textit{(``Did she buy a lot, Mom?'' – ``Just one pastry. But I was happy the whole afternoon.'')}
\end{truyen}
\begin{baihoc}
  Tôn trọng người khác không phụ thuộc vào địa vị của họ. Dành vài phút chân thành hỏi thăm một người xa lạ đôi khi mang lại niềm vui lớn hơn bất kỳ món quà vật chất nào.
\end{baihoc}

\section{Cuộc Tranh Luận Trên Mạng}
\begin{truyen}{Cuộc Tranh Luận Trên Mạng}{Kiềm Chế}
\chuhoa{A}{ }nh viết xong bình luận dài gay gắt phản bác ý kiến của người lạ trên mạng xã hội.\\
\textit{(He finished writing a long, sharp comment refuting a stranger's opinion on social media.)}

Ngón tay chuẩn bị nhấn ``Gửi''.\\
\textit{(His finger was about to press ``Send''.)}

Anh dừng lại, đọc lại một lần. Rồi xóa đi.\\
\textit{(He paused, read it over once more. Then deleted it.)}

Thay vào đó anh nhắn: ``Tôi có quan điểm khác. Nếu bạn muốn trao đổi thêm thì tôi sẵn lòng.''\\
\textit{(Instead he wrote: ``I have a different view. If you would like to discuss, I am open to it.'')}

``\textbf{Trên mạng, từ khó nói nhất không phải là lời chửi mắng. Mà là câu mời ngồi xuống cùng nhau}'' – anh nghĩ thầm.\\
\textit{(``\textbf{Online, the hardest words to type are not insults. They are an invitation to sit down together}'' – he thought to himself.)}
\end{truyen}
\begin{baihoc}
  Kiềm chế trong tranh luận không phải yếu đuối mà là khôn ngoan. Cách ta đối thoại với người bất đồng quan điểm phản ánh trình độ văn hóa và nhân cách của ta rõ ràng hơn bất kỳ ý kiến nào ta phát biểu.
\end{baihoc}

\section{Ghế Ưu Tiên}
\begin{truyen}{Ghế Ưu Tiên}{Tức Thì}
\chuhoa{T}{ }rên xe buýt đông, một người đang ngủ gật trên ghế ưu tiên khi người khuyết tật bước lên.\\
\textit{(On a crowded bus, a man was dozing in a priority seat when a person with a disability boarded.)}

Không ai lên tiếng vì không muốn đụng chạm.\\
\textit{(No one spoke up, unwilling to cause conflict.)}

Một cậu học sinh lớp tám chạm nhẹ vào vai người đang ngủ: ``Chú ơi, có người cần ghế này.''\\
\textit{(A grade eight student gently tapped the sleeping man's shoulder: ``Sir, someone needs this seat.'')}

Người đàn ông thức giấc ngơ ngác rồi đứng dậy ngay, xin lỗi.\\
\textit{(The man woke up in confusion then stood up immediately, apologising.)}

Cậu học sinh không cảm thấy mình làm gì đặc biệt. Nhưng \textbf{cả xe buýt hôm đó học được thứ gì đó từ một đứa trẻ}.\\
\textit{(The student did not feel he had done anything special. But \textbf{the entire bus that day learned something from a child}.)}
\end{truyen}
\begin{baihoc}
  Đối nhân xử thế tốt đôi khi chỉ là thay mặt cộng đồng nói điều đúng đắn vào đúng thời điểm, dù người nói đó chỉ là một đứa trẻ. Tuổi tác không quyết định phẩm cách.
\end{baihoc}

\section{Người Ghi Tên}
\begin{truyen}{Người Ghi Tên}{Nhớ Đến Người Khác}
\chuhoa{T}{ }ại buổi nhận giải thưởng lớn, diễn giả liệt kê rất nhiều người đã giúp mình.\\
\textit{(At a major award ceremony, the speaker listed many people who had helped her.)}

Trong số đó có tên người lái xe đưa đón bà từ hội nghị đến hội nghị suốt ba năm qua.\\
\textit{(Among them was the driver who had ferried her between conferences for the past three years.)}

Người dẫn chương trình ngạc nhiên: ``Bà cũng cảm ơn cả lái xe ạ?''\\
\textit{(The host was surprised: ``You are also thanking the driver?'')}

``Ông ấy giúp tôi đến đúng giờ để làm công việc của mình. Không có ông ấy, tôi không có mặt ở đây. \textbf{Không ai thành công một mình}.''\\
\textit{(``He helped me arrive on time to do my work. Without him, I would not be here. \textbf{No one succeeds alone}.'')}
\end{truyen}
\begin{baihoc}
  Người biết ơn những đóng góp dù nhỏ bé của người khác là người đã học được bài học lớn nhất của cuộc đời: mọi thành công đều được xây trên vai của nhiều người vô danh.
\end{baihoc}

\section{Người Lạ Trong Bệnh Viện}
\begin{truyen}{Người Lạ Trong Bệnh Viện}{Đồng Cảm}
\chuhoa{Ô}{ }ng ngồi một mình trong hành lang bệnh viện, mặt tái nhợt lo lắng.\\
\textit{(He sat alone in the hospital corridor, looking pale and deeply worried.)}

Người phụ nữ đi qua dừng lại: ``Anh ổn không?''\\
\textit{(A woman passing by stopped: ``Are you okay?'')}

``Vợ tôi đang trong phòng mổ. Đã bốn tiếng.''\\
\textit{(``My wife is in surgery. It has been four hours.'')}

Bà ngồi xuống bên cạnh ông, không nói gì thêm. Chỉ ở đó.\\
\textit{(She sat down beside him, saying nothing more. Just being there.)}

Một tiếng sau bác sĩ ra báo tin tốt. Ông quay sang cảm ơn bà. Bà mỉm cười: ``\textbf{Tôi chỉ ngồi thôi. Nhưng tôi nghĩ người đang chờ một mình không nên ngồi một mình}.''\\
\textit{(An hour later the doctor emerged with good news. He turned to thank her. She smiled: ``\textbf{I only sat. But I thought a person waiting alone should not be alone}.'')}
\end{truyen}
\begin{baihoc}
  Đồng cảm không cần lời nói hay hành động đặc biệt. Đôi khi chỉ cần chọn không bỏ đi khi một người lạ đang cần ai đó ở cạnh là đủ để thay đổi cả một khoảnh khắc cuộc đời họ.
\end{baihoc}
